\chapter{Анализ существующих методов} \label{ch1}

В данной главе рассматриваются основные подходы к приоритизации тестовых случаев, их преимущества и ограничения. 
В параграфе~\ref{ch1:sec1} описана задача регрессионного тестирования и приоритизации тестовых случаев.
В параграфе~\ref{ch1:sec2} рассмотрены существующие методы приоритизации и их ограничения.

\section{Регрессионное тестирование и приоритизация тестовых случаев} 
\label{ch1:sec1}

Регрессионное тестирование --- процесс повторного выполнения тестов после внесения изменений в программный код с 
целью подтверждения того, что ранее работавшая функциональность не нарушена. По мере развития проекта объём 
тестового набора растёт, что делает полный повторный прогон нецелесообразным с точки зрения временных и 
вычислительных затрат.

Для решения данной проблемы используются три основных подхода: минимизация тестового набора (Test Suite Minimization), 
отбор тестов (Test Case Selection) и приоритизация тестовых случаев (Test Case Prioritization, TCP). 
В отличие от первых двух, TCP не исключает тесты из набора, а только упорядочивает их выполнение по некоторому 
критерию, сохраняя полноту покрытия.

\section{Существующие методы приоритизации}\label{ch1:sec2}

Среди методов приоритизации выделяют несколько классов, каждый из которых имеет характерные ограничения.

Случайный поиск (Random Search) упорядочивает тесты в произвольном порядке без учёта каких-либо критериев. 
Несмотря на простоту, он служит нижней границей сравнения --- любой разумный алгоритм приоритизации должен 
превосходить случайный порядок.

Жадные алгоритмы на основе покрытия кода упорядочивают тесты по убыванию числа новых покрываемых элементов. 
Алгоритм дополнительного покрытия (Additional) на каждом шаге выбирает тест с максимальным числом ещё не покрытых
элементов --- это простой и эффективный подход, однако он не учитывает вероятность дефектов и принимает локально 
оптимальные решения, не гарантируя глобального оптимума.

Эволюционные методы, в частности генетические алгоритмы, рассматривают задачу TCP как задачу оптимизации перестановок.
В работе~\cite{paygude2020} рассматривается применение ГА к набору данных BigFaultMatrix, показано превосходство 
над случайным поиском и алгоритмами hill climbing. Вместе с тем фитнесс-функция учитывает только скорость 
обнаружения дефектов, не принимая во внимание вероятность дефектов и структурное разнообразие тестов.

В статье~\cite{mahdieh2021} рассматривается подход CovClustering, в котором вероятность дефектов (fault-proneness)
методов оценивается на основе анализа влияния изменений (Change Impact Analysis, CIA). Подход учитывает вероятность
дефектов, однако использует кластеризацию вместо эволюционного поиска, что ограничивает возможности глобальной 
оптимизации порядка тестов.

Таким образом, ни один из рассмотренных методов не сочетает одновременно эволюционный поиск, учёт вероятности 
дефектов и структурное разнообразие тестов. Устранение этих ограничений является целью настоящей работы.